\subsubsection{LLM Forecasting Methods}


Multiple studies have measured zero-shot LLM forecasting capability (zero-shot means without any complex prompt architectures or examples specific to the forecasting task), and found that better general ability base models tend to perform better on forecasting tasks \citep{halawiApproachingHumanLevelForecasting2024} \citep{kargerForecastBenchDynamicBenchmark2024}: In one study with dozens of base models and a dynamically updating benchmark on prediction market forecasting questions, an inverse linear relationship was found between the human preference of a model's answer (in terms of an ELO score) and the Brier score \citep{kargerForecastBenchDynamicBenchmark2024}. A similarly inverse log-linear relationship was found between the compute used to train the model and the Brier score.
Several LLM forecasting methods were investigated to improve overall rating forecasts and narrative forecasting of outcomes.


\emph{deepseek-V3.2} (with reasoning enabled) was selected as the forecasting model for overall ratings, due to its high performance on general reasoning ability benchmarks and low cost. \emph{gemini-3-pro-preview} was also used, due to its superior performance on reasoning benchmarks. Lastly \emph{gemini-2.5-flash
} was selected only to assess the possibility of fine-tuning, as other more powerful models were unavailable for fine-tuning.
 
The LLM forecasting method was decided upon by iteratively inspecting both the quality of the response, and the overall accuracy of the forecasts made by the LLM. \emph{gemini-2.5-pro} was prompted with a series of ``mock forecasts'', generated by \emph{gemini-2.5-pro}. \emph{gemini-2.5-pro} was selected for its high reasoning ability and prompt adherence, as well as its ability to efficiently process PDF documents.  The ``mock forecast'' used relevant pages retrieved by ranking the categorized topics by forecast informativeness and retrieving the most relevant activity description and evaluation pages.

To generate each mock forecast, a retrieval-augmented input was constructed consisting of up to 10 baseline pages and up to 10 outcome/evaluation pages per activity. Baseline pages were selected from high-scoring passages across 19 of the 23 baseline subcategories (objectives, implementation plans, risks, financing details, contextual challenges, and stakeholder/implementer information; excluding condensed summaries, quantitative targets, budget and legal pages, and monitoring detail pages). A minimum relevance score of 9/10 was required. Manual inspection revealed that below a threshold of 3/10 relevance, pages were typical not relevant to reasons for ratings or explaining outcomes. Therefore a relevance threshold of a minimum score of 3/10 was used for page relevance grades. For pages that did match, nearby adjacent pages were also included when insufficient high-scoring pages were available to prevent relevant information and informative tables from being cut off. Pages were selected from post-activity documents which emphasized deviations from plans (including deviations, delays/early completion, and over/under-spending), using a lower relevance threshold (minimum score of 3) and likewise included surrounding pages to reach the target count when needed. These retrieved excerpts were then merged with activity metadata (title, scope, planned start/end dates, planned financing totals when present) and brief model-generated baseline summaries (activity description and risk summary). Gemini was then prompted to write a forecast from the ex-ante perspective. Importantly, the prompt required the model to end by outputting the \emph{known} final evaluation rating for that activity (derived from the merged ratings file and converted into scale-specific text via \texttt{get\_ratings\_text}). The prompt also instructed the model to ground the narrative in the retrieved evaluation pages and to return ``NO RESPONSE'' if the evaluation excerpts did not contain sufficient justification for the assigned rating.


The most semantically relevant activities which ended approximately at or before the start of the activity being forecasted were then retrieved (see Section \ref{ssub:nearest_neighbor_vector_similarity}). In addition, the activity ``risks'' were inserted before each mock forecast, to provide context for the example. Each mock forecast was structured in a way similar to the highest performing scratchpad method from \citep{halawiApproachingHumanLevelForecasting2024}.
A series of features including the activity title, start date, and activity location were injected into the prompt to provide context for the activity, as well as a \emph{gemini-2.5-flash}-generated summary.

% For Method \#2, the most relevant pages of the activity documents were directly converted to text and inserted into the prompt.

Finally, the distribution of rating outcomes was inserted into the prompt, in order to prevent collapse towards only a few ratings.
% LaTeX has no \subsubsubsection by default
% \newcommand{\subsubsubsection}[1]{\paragraph{#1}\mbox{}\\}

\subsubsubsection{Prompt Template and Multi-stage Prompting}

The full prompt is composed of the following blocks, assembled in the order listed:

\begin{enumerate}[noitemsep,topsep=2pt]
  \item \textbf{Activity metadata} --- title, planned start/end dates, scope, planned disbursement (USD), location(s), GDP per capita, reporting organisation.
  \item \textbf{Baseline description \& targets} --- a \emph{gemini-2.5-flash}-generated activity description, targets summary, contextual challenge description, complexity details, integratedness description, financing summary, and implementer-performance context.
  \item \textbf{Risk summary} --- a \emph{gemini-2.5-flash}-generated structured summary of key project risks.
  \item \textbf{Retrieved baseline excerpts (RAG)} --- additional activity-specific information retrieved by the model from the available documents and synthesised into a short passage (\texttt{\{rag\_synthesis\_additional\_info\}}).
  \item \textbf{KNN few-shot block} --- summarised mock forecasts of the three most semantically similar training activities, one from each broad rating tier (see Section~\ref{ssub:nearest_neighbor_vector_similarity}).
  \item \textbf{Multi-stage reasoning (S1/S2)} --- outputs of two prior LLM calls: reasons the activity may have been rated ``Moderately Satisfactory'' or worse (S1) and reasons it may have been rated ``Satisfactory'' or ``Highly Satisfactory'' (S2). This left approximately 50\% if responses above and below the threshold for consideration.
  \item \textbf{Rating distribution prior} --- a short description of the empirical distribution of rating outcomes in the training data, prepended to the scratchpad method instruction block (not appended as the final block).
\end{enumerate}

The full prompt template for the LLM Forecast is shown in Figure~\ref{fig:forecast-prompt-multistage}. LLMs in all prompts were instructed to ``think like a superforecaster'' \citep{tetlockSuperforecastingArtScience2015}. Superforecasters typically outperform domain experts by maintaining broad judgement and cultivating calibrated probability estimates, rather than relying on a single strategy such as econometric analysis. This replicates prior methods which found this instruction to be the highest performing in the geopolitical domain \cite{halawiApproachingHumanLevelForecasting2024}.

\textbf{Few-Shot Block}
In both methods, a $k$-nearest-neighbors (KNN) few-shot block of semantically similar activities in the training data is used (see Section \ref{ssub:nearest_neighbor_vector_similarity} for how semantic similarity was determined). A range of nearest neighbors was evaluated. The language model was prompted to extrapolate lessons about rating scales from the most similar examples in each rating tier: the most similar ``Highly Unsatisfactory'', ``Unsatisfactory'', or ``Moderately Unsatisfactory'' example; the most similar ``Moderately Satisfactory'' example; and the most similar ``Satisfactory'' or ``Highly Satisfactory'' example. A selection of $k=3$ summarized mock forecasts was found to perform better than $k=1$, $5$, or $7$.

Each example activity in the few-shot block included (i) key metadata (title and, where available, location and a brief summary), (ii) a short ``risks'' summary, (iii) the retrospective mock forecast analysis, and (iv) the final evaluation outcome label.

\textbf{Additional Prompts}
Two additional prompts were given, and inserted into the final forecast: (1) reasons the activity may have been evaluated as ``Moderately Satisfactory'' or worse, (2) reasons the activity may have been evaluated as ``Moderately Satisfactory'' or better.

The forecasting prompt required a structured response format that explicitly considered arguments for both lower and higher ratings, and ended with a single-line prediction. Concretely, the model was instructed to: (1) provide reasons the overall success might be rated \texttt{\{midpoint\_low\_text\}} or lower, (2) provide reasons it might be rated \texttt{\{midpoint\_high\_text\}} or higher, (3) aggregate considerations and select exactly one of the \texttt{\{num\_options\}} outcomes, and (4) output the final forecast on the last line beginning with \texttt{FORECAST: } followed by only the chosen option.

Finally, a short description of the empirical distribution of rating outcomes in the training data was appended. This was found to reduce mode-collapse toward a narrow subset of ratings.


% \begin{figure}[htbp]
%   \centering
%   \begin{subfigure}{0.48\textwidth}
%     \begin{tcolorbox}[left=4pt, right=4pt, top=4pt, bottom=4pt]
% \ttfamily\footnotesize
% SYSTEM:\newline
% You are an experienced international aid decision maker with a quantitative mindset.
% Forecast the overall evaluation rating from the options:
% \{options\_text\}.\newline
% USER:\newline
% Forecast what the outcome will be for this activity.\newline
% \#\#\# EXAMPLE ACTIVITIES \#\#\#\newline
% [For each neighbor: title; risks; example forecast;
% rating scale; final evaluation outcome]\newline
% \#\#\# NEW ACTIVITY TO FORECAST \#\#\#\newline
% ACTIVITY ID: \{activity\_id\}\newline
% ACTIVITY TITLE: \{activity\_title\}\newline
% ORIGINAL PLANNED START DATE: \{planned\_start\}\newline
% ORIGINAL PLANNED END DATE: \{planned\_end\}\newline
% ACTIVITY LOCATION(S): \{location\}\newline
% ACTIVITY DESCRIPTION (SUMMARY): \{gemini\_generated\_summary\}\newline
% ACTIVITY RISKS: \{risks\_summary\}\newline
% Provide the following format for your response:\newline
% 1. Provide reasons why the overall success might be rated \{midpoint\_low\_text\} or lower.\newline
% 2. Provide reasons why the overall success might be rated \{midpoint\_high\_text\} or higher.\newline
% 3. Aggregate your considerations, and decide on the final outcome among the \{num\_options\} options.\newline
% 4. Provide the final forecast on the last line beginning with 'FORECAST: '
% followed by only the forecast with no extra words.\newline
% [Append: training-set rating distribution text]\newline
% Respond only in English.
%     \end{tcolorbox}
%     \caption{Method \#1 (summary injection).}
%   \end{subfigure}\hfill
%   \begin{subfigure}{0.48\textwidth}
%     \begin{tcolorbox}[left=4pt, right=4pt, top=4pt, bottom=4pt]
% \ttfamily\footnotesize
% SYSTEM:\newline
% You are an experienced international aid decision maker with a quantitative mindset.
% Forecast the overall evaluation rating from the options:
% \{options\_text\}.\newline
% USER:\newline
% Forecast what the outcome will be for this activity.\newline
% \#\#\# EXAMPLE ACTIVITIES \#\#\#\newline
% [Same few-shot block as Method \#1]\newline
% \#\#\# NEW ACTIVITY TO FORECAST \#\#\#\newline
% ACTIVITY ID: \{activity\_id\}\newline
% ACTIVITY TITLE: \{activity\_title\}\newline
% ORIGINAL PLANNED START DATE: \{planned\_start\}\newline
% ORIGINAL PLANNED END DATE: \{planned\_end\}\newline
% ACTIVITY LOCATION(S): \{location\}\newline
% EXCERPTS FROM BASELINE ACTIVITY DOCUMENTS:\newline
% \{PDF\_to\_text\_excerpts\}\newline
% ACTIVITY RISKS: \{risks\_summary\}\newline
% Provide the following format for your response:\newline
% 1. Provide reasons why the overall success might be rated \{midpoint\_low\_text\}.\newline
% 2. Provide reasons why the overall success might be rated \{midpoint\_high\_text\}.\newline
% 3. Aggregate your considerations, and decide on the final outcome among the \{num\_options\} options.\newline
% 4. Provide the final forecast on the last line beginning with 'FORECAST: '
% followed by only the forecast with no extra words.\newline
% [Append: training-set rating distribution text]\newline
% Respond only in English.
%     \end{tcolorbox}
%     \caption{Method \#2 (raw text injection).}
%   \end{subfigure}
%   \caption{Prompt templates for LLM forecasting Methods \#1 and \#2. The methods share the same scaffold (few-shot examples, metadata, risks, and response-format constraints) and differ only in how activity context is injected (summary vs.\ raw document text).}
%   \label{fig:forecast-prompts}
% \end{figure}
\newtcolorbox[auto counter, number within=section]{boxedprompt}[2][]{
  enhanced,
  breakable,
  colback=white,
  colframe=black,
  fonttitle=\bfseries,
  title=Box~\thetcbcounter: #2,
  #1
}
\begin{boxedprompt}
[label={box:multistage}]{Single-method multi-stage forecasting prompt}{

Stages s1 and s2 are run as separate calls, and their outputs are inserted into the final (s3) prompt via \texttt{insert\_stage\_s1\_answer\_here} and \texttt{insert\_stage\_s2\_answer\_here}. Below are the system and user prompts.

\vspace{1em}
% Single-method multi-stage forecasting prompt
% Stages s1 and s2 are run as separate calls, and their outputs are inserted into the final (s3) prompt via \{\texttt{insert\_stage\_s1\_answer\_here}\} and \{\texttt{insert\_stage\_s2\_answer\_here}\}.
}
\ttfamily
\fontsize{10}{10}\selectfont
\setlength{\baselineskip}{0.9\baselineskip}
SYSTEM:\newline
You are an experienced international aid decision maker with a quantitative mindset.
Respond with a comprehensive, thorough forecast of what the overall evaluation rating of the activity will be,
from the options of \{options\_text\}.\newline

USER:\newline
Forecast what the outcome will be for this activity.\newline

\#\#\# Lessons from similar activities \#\#\#\newline
\{knn\_summary\_text\}\newline
% (If no summary is available, this block can instead be the raw few-shot examples:) \newline
% \{few\_shot\_example\_activities\}\newline
\#\#\# End lessons \#\#\#\newline

\#\#\# Additional specific information about the activity that you summarized \#\#\#\newline
\{rag\_synthesis\_additional\_info\}\newline
\#\#\# End of additional information you summarized \#\#\#\newline

ACTIVITY ID: \{activity\_id\}\newline
ACTIVITY TITLE: \{activity\_title\}\newline
ORIGINAL PLANNED START DATE: \{planned\_start\}\newline
ORIGINAL PLANNED END DATE: \{planned\_end\}\newline
ACTIVITY SCOPE: \{activity\_scope\}\newline
PLANNED TOTAL DISBURSEMENT (USD): \{planned\_total\_disbursement\_usd\}\newline
ACTIVITY LOCATION(S): \{locations\}\newline
LOCATION GDP PER CAPITA, USD: \{gdp\_percap\}\newline
PARTICIPATING ORGANIZATIONS: \{reporting\_orgs\}\newline
IMPLEMENTING ORGANIZATION CATEGORY: \{either ``Government'' or ``NGO'', otherwise line not inserted\}\newline
